\begin{center}
{{\bf\fontsize{14pt}{14.5pt}\selectfont 
MỞ ĐẦU}}
\end{center}
\addcontentsline{toc}{chapter}{MỞ ĐẦU}



Origami truyền thống trong văn hóa dân gian Nhật Bản vốn được biết đến như nghệ thuật tạo hình từ một tờ giấy vuông phẳng, không dùng kéo hay hồ dán. Tuy nhiên, dưới lăng kính toán học hiện đại nửa cuối thế kỷ 20, mỗi nếp gấp thẳng thực chất là sự hiện thực hóa của một thao tác hình học xác định. Khi gấp một tờ giấy, ta đang thực hiện các phép biến hình đối xứng trục, tìm giao điểm và dựng đường bao.

Trong hơn 2.000 năm, hình học cổ điển bị chi phối bởi các tiên đề Euclid cùng hai công cụ chuẩn tắc: thước kẻ không chia vạch và compa. Hệ công cụ này đã đặt ra hai bài toán dựng hình kinh điển thời Hy Lạp cổ đại:
\begin{enumerate}
    \item \textbf{Chia ba một góc:} Chia một góc phẳng bất kỳ thành ba phần bằng nhau.
    \item \textbf{Nhân đôi khối lập phương (Bài toán Delian):} Dựng cạnh của khối lập phương có thể tích gấp đôi thể tích khối lập phương ban đầu (tương đương dựng đoạn thẳng có độ dài $\sqrt[3]{2}$).
\end{enumerate}

Vào thế kỷ 19, lý thuyết Galois và đại số trừu tượng của Pierre Wantzel (1837) cùng Ferdinand von Lindemann (1882) đã chứng minh dứt khoát rằng cả ba bài toán trên đều bất khả thi nếu chỉ dùng thước kẻ và compa. 

Tuy nhiên, giới hạn đó không thuộc về bản chất của hình học, mà thuộc về sự giới hạn của chính hai công cụ cổ điển. Bằng việc hệ thống hóa các thao tác gấp qua hệ 7 tiên đề Huzita--Hatori, Origami đã chứng minh năng lực vượt trội: giải quyết trọn vẹn bài toán chia ba góc và nhân đôi khối lập phương. 

Tiểu luận này tập trung làm rõ:
\begin{itemize}
    \item Giới hạn đại số của thước kẻ -- compa so với năng lực giải phương trình bậc ba của Origami thông qua Tiên đề 6.
    \item Trực quan hóa hình học bản chất của Tiên đề 6: bài toán dựng tiếp tuyến chung của hai đường parabol.
    \item Hướng dẫn thực nghiệm, đo đạc sai số và phân tích hình học cho hai quy trình gập kinh điển của Abe Hisashi và Peter Messer.
\end{itemize}